\documentclass{article}
\usepackage{mathtools} 
\usepackage{fontspec}
\usepackage[UTF8]{ctex}
\usepackage{amsthm}
\usepackage{mdframed}
\usepackage{xcolor}
\usepackage{amssymb}
\usepackage{amsmath}


% 定义新的带灰色背景的说明环境 zremark
\newmdtheoremenv[
  backgroundcolor=gray!10,
  % 边框与背景一致，边框线会消失
  linecolor=gray!10
]{zremark}{说明}

% 通用矩阵命令: \flexmatrix{矩阵名}{元素符号}{行数}{列数}
\newcommand{\flexmatrix}[4]{
  \[
  #1 = \begin{pmatrix}
    #2_{11}     & #2_{12}     & \cdots & #2_{1#4}   \\
    #2_{21}     & #2_{22}     & \cdots & #2_{2#4}   \\
    \vdots      & \vdots      & \ddots & \vdots     \\
    #2_{#31}    & #2_{#32}    & \cdots & #2_{#3#4}
  \end{pmatrix}
  \]
}

% 简化版命令（默认矩阵名为A，元素符号为a）: \quickmatrix{行数}{列数}
\newcommand{\quickmatrix}[2]{\flexmatrix{A}{a}{#1}{#2}}

\begin{document}
\title{注释 6.2}
\author{张志聪}
\maketitle
\begin{zremark}
  例14 为什么可以使用洛必达法则。
\end{zremark}

\textbf{证明：}

有
\begin{align*}
  f'(0) = \lim\limits_{x \to 0} \frac{g(x)}{x^2}
\end{align*}

因为
\begin{align*}
  g'(0) & = \lim\limits_{x \to 0} \frac{g(x) - g(0)}{x - 0} \\
        & = \lim\limits_{x \to 0} \frac{g(x)}{x}
\end{align*}

于是我们有
\begin{align*}
  \lim \limits_{x \to 0} g(x) & = \lim\limits_{x \to 0} \frac{g(x)}{x} x                       \\
                              & = \lim\limits_{x \to 0} \frac{g(x)}{x} \lim\limits_{x \to 0} x \\
                              & = g'(0) \lim\limits_{x \to 0} x                                \\
                              & = 0 \cdot 0                                                    \\
                              & = 0
\end{align*}
于是，定理6.7的条件(i)满足。

因为$g''(0)$存在，按照导数的定义，在点$0$的某个邻域$U(0)$内$g'(x)$有定义，
于是，定理6.7条件(ii)满足。

\textcolor{red}{注：$g''(x)$在$0$处不一定连续，也就是说，在邻域$U^o(0)$上，$g$的二阶导数不一定存在，
也就不能再次使用洛必达法则了。}

\end{document}